\reviewsection{What I Noticed Before I Had the Language}

Before these readings, I already knew what it felt like to be interpreted through fragments. As a neurodivergent person, I have lived through institutions that noticed missed deadlines, uneven performance, communication differences, or difficulty navigating an inaccessible process before they noticed effort, context, intelligence, or persistence. That history shapes how I now look at students. When a child hesitates, repeats a question, moves away from a group, keeps a familiar device close, or needs more time, I do not see a complete explanation. I see a clue.

That is the most important shift Module 1 helped me name. The world does not encounter developmental difference directly. It encounters behavior, language, stories, reports, family knowledge, diagnostic criteria, media images, school routines, scientific claims, and policy categories. Public-health and general educational resources give the world recognizable definitions of autism \citep{cdcASD,simon2016asd,readingRocketsASD}. Those definitions can be useful, but \citet{orsiniDavidson2013critical} show what is lost when statistical urgency turns varied autistic lives into one easily understood crisis. Their warning validated something I have felt personally: a system can collect accurate fragments and still produce an inaccurate person.

The historical reading deepened that point. \citet{harvey2018story} presents autism as a changing story rather than a timeless object that institutions simply discovered. The categories, explanations, and expectations surrounding autism have changed as medicine, schools, families, researchers, media, and autistic communities have changed. My central learning is therefore personal as well as academic: developmental difference is neurological and real, but the meanings attached to it are historical, relational, cultural, and political. No single system gets to explain the whole person.
